% Part: first-order-logic
% Chapter: axiomatic-deduction
% Section: deduction-theorem-quantifiers

\documentclass[../../../include/open-logic-section]{subfiles}

\begin{document}

\olfileid{fol}{axd}{ddq}

\olsection{مبرهنة الاستنباط مع المكمّمات}

\begin{thm}[مبرهنة الاستنباط]
\ollabel{thm:deduction-thm-q} تكون~$\Gamma \cup \{!A\} \Proves !B$
إذا وفقط إذا كان~$\Gamma \Proves !A \lif !B$.
\end{thm}

\begin{proof}
% استكمال برهاني (0120-I1): أُثبت الاتجاه العكسي، وأُتمّت حالة QR الوجودية التي تركها الأصل تمرينًا.
اتجاه ``إذا'' هو الاتجاه المباشر المثبت في
\olref[ded]{thm:deduction-thm}، ولا يعتمد على صورة آخر قاعدة. ولإثبات
اتجاه ``فقط إذا'' في الاشتقاقات الكمية، نكمل الاستقراء السابق بحالتي~\QR{}.

نسير مرة أخرى بالاستقراء في طول !!{derivation} الصيغة $!B$ من
$\Gamma \cup \{!A\}$.

وبرهان أساس الاستقراء مطابق لما في برهان
\olref[ded]{thm:deduction-thm}.

وفي خطوة الاستقراء، افترض مرة أخرى أن !!{derivation} الصيغة $!B$ من
$\Gamma \cup \{!A\}$ ينتهي بخطوة~$!B$ تبررها قاعدة استدلال. فإذا كانت
قاعدة الاستدلال هي إثبات المقدَّم، سرنا كما في برهان
\olref[ded]{thm:deduction-thm}. وإذا كانت قاعدة الاستدلال هي \QR، علمنا
أن $!B \ident !C \lif \lforall[x][!D(x)]$، وأن !!a{formula} من الصورة
$!C \lif !D(a)$ ترد في ال!!{derivation} قبل ذلك، حيث لا يرد $a$ في
$!C$ ولا في $!A$ ولا في $\Gamma$. ومن ثم لدينا
\begin{align*}
  \Gamma \cup \{!A\} & \Proves !C \lif !D(a),\\
  \intertext{وينطبق فرض الاستقراء، أي إن لدينا}
    \Gamma & \Proves !A \lif (!C \lif !D(a)).\\
  \intertext{ومن}
  & \Proves (!A \lif (!C \lif !D(a))) \lif ((!A \land !C) \lif !D(a))\\
  \intertext{وقاعدة إثبات المقدَّم نحصل على}
  \Gamma & \Proves (!A \land !C) \lif !D(a).\\
  \intertext{وبما أن شرط المتغير المميَّز ما يزال منطبقًا، يمكننا أن نضيف إلى هذا ال!!{derivation} خطوة تبررها \QR، فنحصل على}
    \Gamma & \Proves (!A \land !C) \lif \lforall[x][!D(x)].\\
    \intertext{ولدينا أيضًا}
    & \Proves ((!A \land !C) \lif \lforall[x][!D(x)]) \lif (!A \lif (!C \lif \lforall[x][!D(x)])),\\
    \intertext{ومن ثم، بقاعدة إثبات المقدَّم،}
    \Gamma & \Proves !A \lif (!C \lif \lforall[x][!D(x)]),
\end{align*}
أي إن $\Gamma \Proves !A \lif !B$.

وتبقى صورة~\QR{} الثانية. فإذا كان
$!B \ident \lexists[x][!D(x)] \lif !C$، فقد سبقت في ال!!{derivation}
الصيغة~$!D(a) \lif !C$، حيث لا يرد~$a$ في~$!C$ ولا في~$!A$ ولا في
$\Gamma$. ويعطي فرض الاستقراء
\[
  \Gamma \Proves !A \lif (!D(a) \lif !C).
\]
وبإعادة ترتيب المقدِّمين في هذه القضية الصادقة منطقيًا نحصل على
\[
  \Gamma \Proves !D(a) \lif (!A \lif !C).
\]
وشرط الثابت المميَّز ما يزال متحققًا، فتجيز~\QR{} الاستنتاج
\[
  \Gamma \Proves \lexists[x][!D(x)] \lif (!A \lif !C).
\]
ثم تعيد قضية صادقة منطقيًا ترتيب المقدِّمين مرة أخرى، فنحصل بقاعدة
إثبات المقدَّم على
\[
  \Gamma \Proves !A \lif (\lexists[x][!D(x)] \lif !C),
\]
أي~$\Gamma \Proves !A \lif !B$. وبهاتين الصورتين لـ~\QR{}، مع حالات
الفرض والبديهية وإثبات المقدَّم المعالجة في البرهان السابق، يتم الاستقراء.
\end{proof}

\begin{prob}
  تحقّق تفصيلًا من بقاء شروط الثابت المميَّز في صورتي~\QR{} في برهان
  \olref[fol][axd][ddq]{thm:deduction-thm-q}.
\end{prob}

\end{document}
